\subsection{Методология комплексной оценки систем Text-to-SQL}

Детальный анализ качества генерации SQL-запросов в архитектурах, использующих RAG, требует многоуровневого подхода. В отличие от классических задач генерации текста, здесь критически важна синтаксическая валидность, семантическая эквивалентность и вычислительная эффективность.

\subsubsection{Формальные метрики соответствия кода}

\textbf{Формальные метрики соответствия кода}

\textbf{Метод ESM}

Метод ESM направлен на проверку структурного соответствия между сгенерированным запросом $Q_{pred}$ и эталонным $Q_{gold}$. Данный подход нивелирует различия в написании, не влияющие на логику (например, порядок условий в предложении WHERE).

\begin{itemize}
    \item \textbf{Процесс декомпозиции:} Запрос парсится в абстрактное синтаксическое дерево (AST) или стандартизированный формат (например, библиотекой sqlglot или внутренними парсерами бенчмарка Spider). 
    \item \textbf{Компонентный анализ:} Оценка производится по отдельным блокам: SELECT, WHERE, GROUP BY, ORDER BY, FROM. Для каждого блока вычисляется точность (Precision) и полнота (Recall).
    \item \textbf{F1-Score компонентов:} Итоговый балл ESM часто дополняется оценкой частичного совпадения (Partial Matching), что позволяет дифференцировать модели, совершившие незначительную ошибку в имени столбца, от моделей, полностью неверно определивших структуру запроса.
\end{itemize}

\textbf{Execution Accuracy (EX)}

EX является наиболее строгой метрикой, так как она подтверждает функциональную идентичность запросов. Результат выполнения $R_{pred}=execute(Q_{pred}, D)$ сравнивается с $R_{gold}=execute(Q_{gold}, D)$

\begin{itemize}
    \item Проблематика «пустых множеств»: Если оба запроса возвращают пустой результат из-за отсутствия данных в тестовой БД, EX может быть ложноположительной. Для решения этой проблемы в современных тестах (например, BIRD) используются «интенсивные» базы данных с гарантированным наличием записей под каждый тип условий.
    \item 

\end{itemize}

Учет порядка (Ordering): Если в эталонном запросе отсутствует ORDER BY, при сравнении результатов $R_{pred}$ и $R_{gold}$ наборы строк рассматриваются как неупорядоченные множества.


\subsubsection{Продвинутые метрики и бенчмарки (BIRD и Spider)}

Для магистерского исследования принципиально различать подходы двух доминирующих бенчмарков.

\textbf{Spider (Cross-domain \& Complex):}

\begin{itemize}
    \item Вводит классификацию запросов по 4 уровням сложности на основе количества джойнов (JOIN), вложенных запросов и агрегатных функций.
    \item Основной упор делается на Component Matching, что позволяет исследователю понять, на каких именно конструкциях (например, HAVING) модель теряет точность.
\end{itemize}

\textbf{BIRD-SQL (Big Bench for Large-scale Database):}

\begin{itemize}
    \item Valid Efficiency Score (VES): Инновационная метрика, которая учитывает не только корректность результата, но и время выполнения запроса. Это критично для RAG-систем, работающих с Big Data, где неоптимальный JOIN может привести к таймауту.
    \item External Knowledge: BIRD проверяет способность модели использовать внешние словари или описание предметной области, передаваемые в контексте (prompt), что напрямую коррелирует с эффективностью RAG.
\end{itemize}

\textbf{Оценка эффективности Retrieval в контуре RAG}

В архитектурах Retrieval-Augmented Generation для задач Text-to-SQL качество итоговой генерации детерминировано релевантностью извлеченных метаданных: схем баз данных, описаний атрибутов и примеров кортежей. В связи с этим необходима изолированная оценка компонента извлечения (Retrieval) с использованием следующих метрик:

\begin{itemize}
    \item \textbf{Schema Linking Precision/Recall ($P_{sl}, R_{sl}$):} Отражает точность и полноту извлечения таблиц и столбцов, необходимых для построения эталонного запроса $Q_{gold}$. Математически полнота извлечения (Recall) определяется как:
    $$R_{sl} = \frac{|S_{pred} \cap S_{gold}|}{|S_{gold}|}$$
    где $S_{pred}$ — множество извлеченных элементов схемы, $S_{gold}$ — множество элементов, присутствующих в эталонном запросе. Значение $R_{sl} < 1.0$ свидетельствует о том, что генератор (LLM) априори лишен информации, необходимой для корректной дедукции SQL-кода.

    \item \textbf{Context Relevancy ($CR$):} Отношение количества релевантных элементов схемы в сформированном промпте к общему объему переданного контекста. Высокая доля нерелевантной информации (шума) в контексте коррелирует с ростом вероятности «галлюцинаций», при которых модель оперирует семантически близкими, но технически неверными идентификаторами объектов БД.
\end{itemize}

\textbf{Интегральная оценка (Faithfulness и SQL-валидность)}

Для комплексного анализа системы предлагается использовать интегральный коэффициент, учитывающий не только совпадение с эталоном, но и формальные свойства кода:

\begin{enumerate}
    \item \textbf{Syntactic Validity ($V_{syn}$):} Доля запросов, успешно прошедших синтаксический анализ (parsing) и не содержащих ошибок компиляции.
    \item \textbf{Semantic Alignment:} Качественная оценка соответствия логики сгенерированного запроса $Q_{pred}$ интенту пользователя. Реализуется через парадигму \textit{LLM-as-a-judge} (например, с применением модели GPT-4o), где проверяется соответствие предикатов и агрегатных функций исходной бизнес-логике вопроса.
\end{enumerate}

\begin{table}[h!]
\centering
\caption{Сравнительный анализ методов оценки Text-to-SQL}
\label{table:metrics_comparison}
\begin{tabular}{|p{2cm}|p{2.5cm}|p{3.5cm}|p{4cm}|}
\hline
\textbf{Метрика} & \textbf{Объект оценки} & \textbf{Преимущества} & \textbf{Недостатки} \\ \hline
ESM & Структура кода & Высокая скорость, интерпретируемость & Игнорирует альтернативные синтаксические конструкции \\ \hline
EX & Результат (data) & Максимальная объективность & Риск ложноположительных срабатываний на пустых БД \\ \hline
VES & Производительность & Учет вычислительной эффективности & Чувствительность к аппаратному обеспечению \\ \hline
Recall@K & Retrieval-контекст & Локализация ошибок на этапе поиска & Не гарантирует валидность финального SQL-кода \\ \hline
\end{tabular}
\end{table}

